%Master File:lectures.tex

\lesson{Training HMMs}
\vspace{-2cm}
\begin{center}
 \includegraphics[height=10cm,angle=90]{hmmresult}
\end{center}

\keywords{Baum Welch, HMM Structures}
%%%%%%%%%%%%%%%%%%%%%%% Next Slide %%%%%%%%%%%%%%%%%%%%%%%
\renewcommand{\Outline}{%
\begin{slide}
\section[1]{Outline}

\begin{minipage}{8cm}
  \begin{enumerate}\squeeze
    \outlineitem{Recap}{reca}
    \outlineitem{Baum Welch}{baumwelch}
    \outlineitem{HMM Structure}{hmmst}
    \outlineitem{Applications}{applications}
   \end{enumerate}
\end{minipage}\hfill
\begin{minipage}{15cm}
  \vfill
  \includegraphics[height=12cm,angle=90]{hmmresult}
  \vfill
\end{minipage}
\end{slide}
\addtocounter{outlineitem}{1}
}

\setcounter{outlineitem}{1}
\Outline
\toptarget{firstoutline}
%%%%%%%%%%%%%%%%%%%%%%% Next Slide %%%%%%%%%%%%%%%%%%%%%%%

\begin{slide}
\section{Recap}

\begin{PauseHighLight}
  \begin{itemize}
  \item HMMs model the generation of sequences\pause
  \item They are probabilistic finite state automata describing the
    transition between state\pause
  \item Each state emits a symbol with some probability\pause
  \item We can compute the joint probability $\Prob{\bm{q},\bm{\xi}}$ of
    the states, $\bm{q}$, and sequence $\bm{\xi}$\pause
  \item Using forward and backward variable we can sum over all possible
    paths through the state\pause
  \end{itemize}
\end{PauseHighLight}

\end{slide}

%%%%%%%%%%%%%%%%%%%%%%% Next Slide %%%%%%%%%%%%%%%%%%%%%%%

\begin{slide}
\section{Structure and Parameters}

\begin{PauseHighLight}
  \begin{itemize}
  \item The structure of the HMM is the set of states and the allowed
    transition between states\pause
  \item Given a structure the parameters, $\bm{\theta}$, describes the
    emission probabilities, $e_i(\zeta)$ at each state and the
    transition probabilities, $a_{ij}$, between states\pause
  \item When we understand the process generating the sequence very well
    then we can sometimes specify the structure and parameters
    \textit{ab initio}\pause
  \item Most of the time, we don't have a detailed understanding of the
    process and have to learn the parameters\pause
  \end{itemize}
\end{PauseHighLight}

\end{slide}


%%%%%%%%%%%%%%%%%%%%%%% Next Slide %%%%%%%%%%%%%%%%%%%%%%%
\Outline % Baum Welch
%%%%%%%%%%%%%%%%%%%%%%% Next Slide %%%%%%%%%%%%%%%%%%%%%%%

\begin{slide}
\section{Baum-Welch Training}

\begin{PauseHighLight}
  \begin{itemize}
  \item We can train the parameters of an HMM to maximise the likelihood
    of a set of training sequences\pause
  \item This algorithm developed mainly by Leonard Baum is known as
    the Baum-Welch algorithm\pause
  \item It is an iterative algorithm that finds a \emph{local maximum}
    of the likelihood of a set of training sequences\pause
  \item It is an example of an \emph{expectation maximisation} (EM)
    strategy\pause
  \end{itemize}
\end{PauseHighLight}

\end{slide}

%%%%%%%%%%%%%%%%%%%%%%% Next Slide %%%%%%%%%%%%%%%%%%%%%%%

\begin{slide}
\section{Iterative Improvement}

\begin{PauseHighLight}
  \begin{itemize}
  \item Given a set of sequence
    $\mathcal{D}=\{\bm{\xi}^\mu|\mu=1,\,\ldots,\,P\}$,
    then if we knew the path through the set of states then we can count
    \begin{itemize}
    \item $A_{ij}$ the number of times we transition from state
      $i$ to $j$\pause
    \item $E_i(\zeta)$ the number of times we emit symbol $\zeta$ when
      in state $i$\pause
    \end{itemize} 
  \item To maximise the likelihood of the particular sequences we can set
    \begin{align*}
      e_i(\zeta) &= \frac{E_i(\zeta)}{\sum\limits_{\xi\in\mathcal{A}} E_i(\xi)},&
      a_{ij} &= \frac{A_{ij}}{\sum\limits_{k\in\mathcal{Q}} A_{ik}}\pause
    \end{align*}
  \end{itemize}
\end{PauseHighLight}

\end{slide}

%%%%%%%%%%%%%%%%%%%%%%% Next Slide %%%%%%%%%%%%%%%%%%%%%%%

\begin{slide}
\section{Pseudo Counts}

\begin{PauseHighLight}
  \begin{itemize}
  \item A problem is that if we use the HMM on unseen data then there
    might be emissions or transitions that we've set to zero which we
    want to allow\pause
  \item To solve that we can add \emph{pseudo counts} giving some small
    probability of always emitting a symbol or making a transition
    \begin{align*}
      e_i(\zeta) &= \frac{E_i(\zeta)+ r_i(\zeta)}
        {\sum\limits_{\xi\in\mathcal{A}} \left(E_i(\xi)+r_i(\xi)\right)}, & 
      a_{ij} &= \frac{A_{ij} + r_{ij}}
        {\sum\limits_{k\in\mathcal{Q}} \left(A_{ik}+r_{ik}\right)}\pause
    \end{align*}
  \item Usually $r_i(\zeta)$ and $r_{ik}$ will be some small number\pause
  \end{itemize}
\end{PauseHighLight}

\end{slide}

%%%%%%%%%%%%%%%%%%%%%%% Next Slide %%%%%%%%%%%%%%%%%%%%%%%

\begin{slide}
\section{Chicken and Egg}

\begin{PauseHighLight}
  \begin{itemize}
  \item The second problem we face though is we don't know the path through the
    states that a sequence should take\pause
  \item That path will depend on the emission and transition
    probabilities (which is what we want to determine)\pause
  \item We therefore use an iterative learning process where we
    initialise the transition and emission probabilities at random\pause
  \item We then compute $E_i(\zeta)$ and $A_{ij}$ for those
    parameters\pause
  \item We then update the parameters as described above and iterate,
    nudging the parameters towards a local optimum\pause
  \end{itemize}
\end{PauseHighLight}

\end{slide}

%%%%%%%%%%%%%%%%%%%%%%% Next Slide %%%%%%%%%%%%%%%%%%%%%%%

\begin{slide}
\section[-1]{Computing the Emission}

\begin{PauseHighLight}
  \begin{itemize}
  \item For a given set of parameters $\bm{\theta}$ the probability of
    emitting a symbol $\zeta$ in state $i$ for the set of training
    sequences is
    \begin{align*}
      E_i(\zeta) = \sum_{\bm{\xi}^\mu\in\mathcal{D}} \sum_{t=1}^{T^\mu}
      \Prob{q_t=i| \bm{\xi}^\mu, \bm{\theta}} \, \pred{\xi_t^\mu = \zeta}\pause
    \end{align*}
  \item Where $\pred{\text{predicate}}$ equals 1 if predicate is true
    and 0 otherwise\pause
  \item But from the backward and forward variables
    \begin{align*}
      E_i(\zeta) = \sum_{\bm{\xi}^\mu\in\mathcal{D}} 
      \frac{1}{\Prob{\bm{\xi}^\mu|\bm{\theta}}}\sum_{t=1}^{T^\mu}
      \alpha_t^\mu(i)\,\beta_t^\mu(i)\, \pred{\xi_t^\mu = \zeta}\pause
    \end{align*}
  \end{itemize}
\end{PauseHighLight}

\end{slide}


%%%%%%%%%%%%%%%%%%%%%%% Next Slide %%%%%%%%%%%%%%%%%%%%%%%

\begin{slide}
\section{Computing Transitions}
\begin{PauseHighLight}
  \begin{itemize}
  \item Similarly the expected number of transitions is given by
    \begin{align*}
      A_{ij} = \sum_{\bm{\xi}^\mu\in\mathcal{D}} \sum_{t=1}^{T^\mu-1}
      \Prob{q_t=i,\, q_{t+1}=j| \bm{\xi}^\mu, \bm{\theta}}\pause
    \end{align*}
  \item After a bit of manipulations we find
    \begin{align*}
      A_{ij} =  \sum_{\bm{\xi}^\mu\in\mathcal{D}} 
      \frac{1}{\Prob{\bm{\xi}^\mu|\bm{\theta}}}\sum_{t=1}^{T^\mu-1}
      \alpha^\mu_t(i)\, a_{ij}\, e_j(\xi_{t+1}) \, \beta^\mu_{t+1}(j)\pause
    \end{align*}
  \item We have already seen that we can compute
    $\Prob{\bm{\xi}^\mu|\bm{\theta}}$ from the forward variable\pause
  \end{itemize}
\end{PauseHighLight}

\end{slide}

%%%%%%%%%%%%%%%%%%%%%%% Next Slide %%%%%%%%%%%%%%%%%%%%%%%

\begin{slide}
\section{Training in Practice}

\begin{PauseHighLight}
  \begin{itemize}
  \item Baum-Welch is widely used\pause
  \item The implementation is slightly non-trivial (must prevent numbers
    underflowing), use sensible pseudo counts\pause
  \item Most practitioners use packages (in python, R, C++,
    matlab)\pause
  \item Also possible to use gradient based algorithms for parameter
    optimisation\pause
  \item More complicated, but better convergence\pause
  \end{itemize}
\end{PauseHighLight}

\end{slide}


%%%%%%%%%%%%%%%%%%%%%%% Next Slide %%%%%%%%%%%%%%%%%%%%%%%
\Outline % HMM Structure
%%%%%%%%%%%%%%%%%%%%%%% Next Slide %%%%%%%%%%%%%%%%%%%%%%%

\begin{slide}
\section{HMM Structures}

\begin{PauseHighLight}
  \begin{itemize}
  \item Whereas training the parameters of an HMM is relatively
    straightforward designing the structure on an HMM is a black
    art\pause
  \item We have to determine
    \begin{itemize}
    \item the number of states\pause
    \item which transition we allow\pause
    \end{itemize}
  \item Sometimes we understand the problem enough to know what we are
    modelling, but this is rather unusual\pause
  \end{itemize}
\end{PauseHighLight}

\end{slide}

%%%%%%%%%%%%%%%%%%%%%%% Next Slide %%%%%%%%%%%%%%%%%%%%%%%

\begin{slide}
\section[-2]{Number of States}

\begin{PauseHighLight}
  \begin{itemize}
  \item A state represents a part of sequence that has a different set
    of emission frequencies\pause
  \item This should be a consequence of some physical differences
    \begin{itemize}
    \item Coding area
    \item Promoter region
    \item Junk DNA\pause
    \end{itemize}
  \item If we have too few states we can fail to model significant
    differences\pause
  \item If we have too many states then we learn chance fluctuations in
    emission frequencies giving poor generalisation\pause
  \end{itemize}
\end{PauseHighLight}

\end{slide}

%%%%%%%%%%%%%%%%%%%%%%% Next Slide %%%%%%%%%%%%%%%%%%%%%%%

\begin{slide}
\section{Transitions}

\begin{PauseHighLight}
  \begin{itemize}
  \item The transitions should represent real physical structure
    \begin{itemize}
    \item Promoter regions precede coding regions\pause
    \end{itemize}
  \item If we have a fully connected HMM we will get the best
    likelihood, but will tend to overfit\pause
  \item Typically experts build quite structured network\pause
  \end{itemize}
\end{PauseHighLight}

\end{slide}

%%%%%%%%%%%%%%%%%%%%%%% Next Slide %%%%%%%%%%%%%%%%%%%%%%%

\begin{slide}
\section{High-Order HMM}

\begin{PauseHighLight}
  \begin{itemize}
  \item The Markov property of HMMs often imposes too little short range
    structure\pause
  \item We can generalise a Markov Chain to second order
    \begin{align*}
      \Prob{x_i|x_{i-1},\,x_{i-2},\,\ldots,x_1} 
      = \Prob{x_i| x_{i-1}, x_{i-2}}\pause
    \end{align*}
  \item Can be further generalised to $n^{th}$-order\pause
  \item In HMM's we often make the emission probabilities dependent on
    several previous elements in a sequence\pause
  \end{itemize}
\end{PauseHighLight}

\end{slide}

%%%%%%%%%%%%%%%%%%%%%%% Next Slide %%%%%%%%%%%%%%%%%%%%%%%

\begin{slide}
\section[-2]{Common Motifs}

\begin{PauseHighLight}
  \begin{itemize}
  \item In biology it is common to build an HMM from building
    blocks\pause
    \begin{center}
      \includegraphics[width=0.4\linewidth]{blocks_1}\\
      \includegraphics[width=0.4\linewidth]{blocks_2}\\
      \includegraphics[width=0.4\linewidth]{blocks_3}
    \end{center}
  \item Use these to model promoter regions\pause
    \begin{center}
      \includegraphics[width=0.9\linewidth]{promoter8}
    \end{center}
  \end{itemize}
\end{PauseHighLight}

\end{slide}

%%%%%%%%%%%%%%%%%%%%%%% Next Slide %%%%%%%%%%%%%%%%%%%%%%%

\begin{slide}
\section{Evolving HMMs}

\begin{PauseHighLight}
  \begin{itemize}
  \item Traditionally HMMs are designed by ``experts''\pause
  \item However by combining common motifs genetic algorithms can also
    be used\pause 
  \item For large problems the GA generated HMMs are actually
    significantly better than the ``expert'' HMMs\pause
  \item It is probably rather hopeful to expect that a human can really
    design an HMM structure with tens of states\pause
  \end{itemize}
\end{PauseHighLight}

\end{slide}

%%%%%%%%%%%%%%%%%%%%%%% Next Slide %%%%%%%%%%%%%%%%%%%%%%%
\Outline % HMM Structure
%%%%%%%%%%%%%%%%%%%%%%% Next Slide %%%%%%%%%%%%%%%%%%%%%%%

\begin{slide}
\section[-2]{DNA Methylation}
\begin{PauseHighLight}
  \begin{itemize}
  \item Most biologists aren't that interested in dishonest
    casinos, but they are interested in promoter regions\pause
    \begin{center}
      \includegraphics[width=0.9\linewidth]{DNAmeth}
    \end{center}
  \end{itemize}
\end{PauseHighLight}

\end{slide}

%%%%%%%%%%%%%%%%%%%%%%% Next Slide %%%%%%%%%%%%%%%%%%%%%%%

\begin{slide}
\section{CpG Island}
\begin{PauseHighLight}
  \begin{itemize}
  \item In mammals non-methylated regions are associate with regions
    with an above average concentration of C followed by a G (known as CpG
    islands)\pause
  \item These are usually under-represented because they are chemically
    modified by methylation (there is a relatively high chance of this
    methyl-C mutating to a T)\pause
  \item However, this methylation is suppressed in some regions of the
    genome (e.g. in promoter regions)\pause
  \item Therefore, finding CpG islands is of biological interest\pause
  \end{itemize}
\end{PauseHighLight}


\end{slide}


%%%%%%%%%%%%%%%%%%%%%%% Next Slide %%%%%%%%%%%%%%%%%%%%%%%

\begin{slide}
\section{Family Resemblance}

\begin{PauseHighLight}
  \begin{itemize}
  \item In the earlier lectures we said that one of the most important
    tasks was to identify similar proteins\pause
  \item However, pairwise comparisons of proteins from distance
    relatives can give rather inconclusive results\pause
  \item However, if we can compare against a whole family of proteins we
    often can spot similarities that are much less obvious\pause
  \item A profile is an alignment of a whole family of proteins\pause
  \item We can train a \emph{profile HMM} to spot a distance relative\pause
  \end{itemize}
\end{PauseHighLight}

\end{slide}

%%%%%%%%%%%%%%%%%%%%%%% Next Slide %%%%%%%%%%%%%%%%%%%%%%%

\begin{slide}
\section{Profile HMMs}

\begin{center}
  \includegraphics[width=0.9\linewidth]{profile}
\end{center}
\end{slide}

%%%%%%%%%%%%%%%%%%%%%%% Next Slide %%%%%%%%%%%%%%%%%%%%%%%

\begin{slide}
\section[-2]{Gene Detection}

\begin{PauseHighLight}
  \begin{itemize}
  \item Gene detection is a significant challenge\pause
  \item The first gene detectors (GENSCAN 1997) found only $\sim10\%$ of
    genes in the human genome with huge number of false positive\pause
  \item Over the years this has improved with new algorithms finding
    around 20\,000 genes in the human genome\pause
  \item Most algorithms use HMMs\pause
    \begin{center}
      \includegraphics[width=0.6\linewidth]{sequenceLogo}
    \end{center}
  \end{itemize}
\end{PauseHighLight}

\end{slide}



%%%%%%%%%%%%%%%%%%%%%%% Next Slide %%%%%%%%%%%%%%%%%%%%%%%

\begin{slide}
\section{Summary}

\begin{PauseHighLight}
  \begin{itemize}
  \item To use HMMs in situations where we don't understand the
    generation process fully we train the parameters using
    Baum-Welch\pause
  \item We still need to determine the architecture of HMMs\pause
  \item Sometimes this is determined by the problem (e.g. profile
    HMM)\pause
  \item Frequently we have to guess and experiment\pause
  \item We can even use optimisation algorithms (although its not that
    easy)\pause
  \end{itemize}
\end{PauseHighLight}

\end{slide}
